% ============================================================================
% Chapter 28 --- Cross-cutting concerns
% ----------------------------------------------------------------------------
% Purpose : the organisational and process concerns that determine whether
%           the playbook actually works in practice. Not a wave; addresses
%           topics that span every wave.
% Page target: ~5 pages.
% Dependencies: Ch. 17--27 (the playbook); Ch. 3 (sovereignty grid for the
%               financial dimension).
% ============================================================================

\chapter{Cross-cutting concerns}
\label{ch:cross-cutting}

The waves of the playbook describe \emph{what} the institution
does, in what sequence, against which criteria. The cross-cutting
concerns of this chapter describe \emph{how} the institution does
it: the organisational disciplines, the communication patterns,
the financial framework, the procurement adjustments without which
the technical work of the waves does not produce the institutional
outcome the playbook intends. They are presented as topical
guidance rather than as prescriptions, since each institution's
existing arrangements differ substantially.

%-----------------------------------------------------------------------------
\section{Change management}
\label{sec:cc-change}

The architecture's transformation interacts with the institution's
existing change-management process. Three patterns recur and merit
explicit handling.

\paragraph{Per-wave change-management cadence.} Each wave produces
a stream of changes that the institution's change advisory board
(or equivalent) must review. The stream is bounded by the wave's
scope and predictable in volume; the cadence agreed at wave open
ensures that the engineering team's pace and the change-management
process's pace are aligned. Misalignment in either direction --- a
backlog of unreviewed changes, or a change-management process
calibrated for a slower pace --- typically produces the wave's
most visible operational difficulties.

\paragraph{ITIL or COBIT integration.} For institutions whose
service-management framework is built on ITIL or COBIT, the
playbook's vocabulary maps onto familiar concepts: a wave is a
change programme, an exit criterion is a change closure, a
reversibility window is a back-out window. The architectural
posture is to use the institution's existing vocabulary where it
fits, rather than imposing the playbook's vocabulary as a
parallel.

\paragraph{Change windows during waves.} Several waves
(particularly Wave~3 network segmentation, Wave~6 endpoint, and
Wave~9 lifecycle) require scheduled disruption of services that
the institution's user constituencies tolerate at agreed
windows. Coordinating these windows across faculties, with
academic-calendar awareness, is a persistent organisational
effort that the playbook does not eliminate.

%-----------------------------------------------------------------------------
\section{Communication}
\label{sec:cc-communication}

A transformation programme of the playbook's scope requires
communication at three levels, each with its own audience,
cadence, and tone.

\paragraph{Governance communication.} The programme governance
body (established in Wave~0) receives structured updates at its
regular cadence: progress against the wave plan, conformance
results from \trchap{ch:conformance}{}, sovereignty grid drift
since the last review, and risk register updates. The governance
update is concise and decision-oriented; it is not a technical
report.

\paragraph{Faculty and operational communication.} Faculties,
departments, and research projects receive communication at the
events that affect them: the cut-over schedule for their building
or service, the changes to their endpoint posture, the
introduction of self-service for their consumers. The pattern is
\emph{advance notice plus contact point}: the affected
constituency is informed before the change, with a named person
to whom they can address concerns.

\paragraph{Institutional and public communication.} The
sovereignty position itself --- the institution's commitment to
the grid, the migration trajectory it has chosen, the precedents
it draws on --- is suitable for communication beyond the IT
audience. The phase-2 international peer comparison anticipated
by this Reference is one of the contexts in which institutional
communication of sovereignty positioning becomes valuable.

%-----------------------------------------------------------------------------
\section{Training and skills development}
\label{sec:cc-training}

The architecture's operation requires skills that may not be
fully present in the institution at the playbook's start. Three
categories of skill investment merit explicit planning.

\paragraph{Per-layer skill development.} Each layer of Part~II
implies specific skills (policy-language authoring for the policy
layer, distributed-storage operation for the storage layer,
container-orchestration administration for the compute layer,
policy-driven infrastructure composition for the orchestration
layer). The institution's training plan addresses these skills
proportionately to its chosen trajectory: a brownfield trajectory
anchored on an incumbent commercial platform invests less in
distributed storage and more in the bridging logic between that
platform and the open policy layer; a brownfield-hybrid trajectory
invests more uniformly across the layers.

\paragraph{Research-network and peer-institution mutualisation.}
Several skills the playbook requires are scarce in any single
institution. The institution's relationship with its national
research and education network, with peer institutions through
its regional research-network federation, and
with any cross-institutional academic network it joins, are
opportunities for shared
training, shared on-call, and shared expertise that no single
institution would sustain alone.

\paragraph{Long-horizon skill investment.} The skills the
playbook builds are not transitional. The institution that
operates the architecture five years after Wave~9 needs the same
skill mix it had at Wave~9. The training investment is
consequently a continuing operational expense rather than a
project expense.

%-----------------------------------------------------------------------------
\section{Risk register}
\label{sec:cc-risk}

A wave-aware risk register is part of the programme's deliverables
from Wave~0 forward. The register's structure follows the
institution's existing risk methodology where one is in place
(\loc{national-risk-method}, ISO/IEC~27005, FAIR, or comparable); the recommendation
of this chapter is that the register be \emph{wave-tagged} so that
each risk is associated with the wave or waves in which it is most
acute.

\paragraph{Per-wave risk profile.} Each wave chapter implicitly
identifies risks: identity-migration risk (Wave~1), policy
miscalibration risk (Wave~4), classification-drift risk (Wave~8).
The register makes these risks explicit and tracks them through
the wave's life. A risk that materialises is documented; a risk
that does not materialise is documented as such at wave closure.

\paragraph{Cross-wave risks.} Some risks span multiple waves: the
risk that the engineering team's capacity is exceeded by parallel
waves; the risk that an institutional change (leadership
transition, funding shift) disrupts the programme's continuity;
the risk that supplier consolidation alters the financial
projections of the chosen trajectory. These risks have no
single-wave home; they are reviewed by the programme governance
at every cycle.

\paragraph{Risk acceptance.} Some risks are accepted explicitly
rather than mitigated. The acceptance is documented, with the
sovereignty-grid dimension on which the risk falls, the
reasoning, and the review schedule.

%-----------------------------------------------------------------------------
\section{Financial planning}
\label{sec:cc-financial}

The playbook implies a financial trajectory that the institution
plans against. Three observations recur across institutional
contexts.

\paragraph{CAPEX-to-OPEX shift.} The reference architecture
typically shifts a portion of the institution's IT spending from
capital expenditure (large infrastructure refreshes) to operating
expenditure (continuous engineering, supplier subscriptions,
training). The shift is not absolute --- substantial CAPEX
remains for compute, storage and network refreshes --- but the
proportion changes in a manner the institution's financial
planning should anticipate.

\paragraph{Sovereignty reinvestment.} A portion of the operating
expenditure that earlier flowed to a single dominant supplier
(licence renewals, premium support, captive services) is
reinvested in institutional capability: engineer salaries, peer
participation, training, contribution to open-source projects on
which the architecture rests. The reinvestment is the financial
form of sovereignty: the same money produces institutional
capability rather than supplier revenue.

\paragraph{Multi-year horizon.} The 5/7/10-year pace options of
\trchap{ch:playbook-reading}{} imply matching financial
horizons. An institution whose financial planning operates on a
single-year cycle will struggle to commit to the playbook
durably. The wave's incremental structure mitigates this --- each
wave can be financially scoped --- but the cumulative trajectory
benefits from a multi-year envelope.

%-----------------------------------------------------------------------------
\section{Procurement adjustments}
\label{sec:cc-procurement}

The sovereignty grid translates into procurement clauses. The
institution's existing \gls{rfp} templates can absorb the
playbook's requirements through targeted additions rather than
through wholesale rewriting.

\paragraph{Standard-interface clauses.} For every procurement
that touches a layer of Part~II, the \gls{rfp} requires the
candidate to demonstrate conformance with the relevant standard
interface (OIDC for identity, OPA REST for policy, NETCONF/YANG
for network, S3 API for object storage, OpenTelemetry for
observability, and so on). The conformance is verified during
the procurement evaluation against the test categories of
\trchap{ch:conformance}{}.

\paragraph{Exit-assistance clauses.} The contract includes a
documented exit-assistance procedure: data export in agreed
formats, configuration export in re-deployable form, support
during a transition period whose duration is set in advance.
The exit cost is bounded contractually and is proportionate to
the genuine migration work rather than to the supplier's
commercial preference.

\paragraph{Sovereignty-gating language.} The institution may
adopt explicit sovereignty gates for specified workload classes:
\emph{the data-dimension score must be Established or better for
sensitive workloads}, \emph{the financial dimension must be
Robust or better for institution-critical components}, and so
on. The gates are declared before the candidate is scored, not
after, and they apply uniformly to commercial and open
candidates alike.

\paragraph{Multi-supplier capability.} For each layer, the
institution preserves the realistic possibility of substituting
the supplier within an agreed window. The procurement framework
ensures that the practical exercise of this possibility is not
prevented by contractual terms (no exclusivity that would
prevent parallel evaluation, no licence terms that would
penalise reduced consumption during transition, no data formats
that would prevent independent re-import).
